\chapter{DENEYSEL DOĞRULAMA}
\label{chap:deney}

Bölüm~\ref{chap:anahtarlama} ve Bölüm~\ref{chap:sin3} içinde benzetim düzeyinde incelenen klasik anahtarlama, nötr anahtar susturma ve susturma+Sin\textsuperscript{3} yöntemleri, geliştirilen üç fazlı T-tipi prototipin tek faz kolunda, düşük güç koşulunda karşılaştırılmıştır. Bölümde deney düzeneği, güç donanımı, denetim ve ölçüm zinciri testleri yeniden üretecek ayrıntıda tanıtılmakta, deney kayıtları değerlendirilmekte ve benzetim modelinin deney özetine kalibrasyonu (Alt Bölüm~\ref{sec:kalibrasyon}) verilmektedir. Deneyler güvenli laboratuvar koşulunda, yaklaşık \SI{2.3}{\ampere} faz akımındaki düşük güç noktasında yürütülmüştür. Anma gücündeki geri beslemeli (S4) çalışma kalibre benzetim modeliyle, GT-tipi şarj yönü ise Bölüm~\ref{chap:gttipi} içinde benzetim kayıtlarıyla incelenmiştir. Deney kayıtlarındaki metin, tablo ve ekran değerleri birbiriyle tutarlılık bakımından denetlenmiştir.

\textbf{Deneysel veri kümesi.} Deneysel veri kümesi, üç yönteme ait PC yazılımı ekranları, osiloskop ekranları ve her yöntem için bir satırlık deney özet tablosundan oluşmaktadır (Tablo~\ref{tab:gelen-excel-senaryolar}). Ölçüm belirsizliği üretici toleransları ve okuma çözünürlüğünden kestirilmiştir. Benzetim modeli deney özetinin aktif ve reaktif gücüne kalibre edilmiş, diğer büyüklükler bağımsız karşılaştırma olarak raporlanmıştır (Alt Bölüm~\ref{sec:kalibrasyon}). Donanım fotoğrafları, devre kartı çizimleri, osiloskop ve PC yazılımı ekranları deney kayıtlarındaki özgün görüntü dosyalarından alınmış ve özgün çözünürlükleriyle verilmiştir. Bölümdeki sayısal değerler deney özet tablosundan ve ekranlarda yazılı değerlerden okunmuştur; ekran üzerinden yapılan tek ölçüm, Alt Bölüm~\ref{sec:thd-kaynaklar} içindeki faz akımı izlerinin sayısallaştırılmasıdır.

\section{Deney Düzeneği ve Güç Donanımı}

Deney düzeneği, yaklaşık 20--25~kW güç dönüştürme kapasitesi hedeflenerek düşey bir pano üzerine bakır baralarla modüler biçimde kurulmuştur. Akım sensörü boyutlandırma hesabında prototip gücü \SI{20}{\kilo\volt\ampere} olarak verilmiştir. Tezde \SI{20}{\kilo\volt\ampere} anma görünür gücü, \num{20}--\SI{25}{\kilo\watt} ise tasarım güç aralığı olarak kullanılmaktadır (Tablo~\ref{tab:sistem-parametreleri}). Şebeke gerilimi güvenlik amacıyla varyak üzerinden yaklaşık \SI{120}{\volt} ile sınırlandırılmış, deneyler yaklaşık \SI{2.3}{\ampere} faz akımında, düşük güç noktasında yapılmıştır.

Şekil~\ref{fig:deney-duzenegi-genel}; varyak modülü, evirici panosu, bilgisayar ve osiloskopları aynı görünümde vermektedir. Varyak modülünde enerji bağlantıları, şalt ve koruma ekipmanları, akım transformatörleri ve AA enerji analizörü bulunmakta; üç faz bağımsız varyaklarla ayarlanabilmektedir. Düzenekte aşırı akım, kaçak akım, topraklama ve acil durdurma korumaları bulunmaktadır.

\begin{figure}[H]
  \centering
  \includegraphics[width=0.95\textwidth]{deneysel/onceki_tez_sekil_5_1_deney_duzenegi.png}
  \caption{Deney ve ölçme düzeneklerinin genel görünümü.}
  \label{fig:deney-duzenegi-genel}
\end{figure}

Şekil~\ref{fig:deney-evirici-genel}, evirici donanımının önden genel görünümünü sunmaktadır. Üç faza ait güç modülleri, kontrol ve ölçüm kartları, DA bara bağlantıları ve yardımcı beslemeler aynı pano üzerindedir; genel görünüm, numaralı ayrıntıların verildiği Şekil~\ref{fig:deney-donanim-kartlar} için yerleşim bağlamı sağlamaktadır.

\begin{figure}[H]
  \centering
  \includegraphics[width=0.88\textwidth]{deneysel/onceki_tez_sekil_5_2_evirici_donanimi.jpeg}
  \caption{Evirici donanımının genel görünümü.}
  \label{fig:deney-evirici-genel}
\end{figure}

Şekil~\ref{fig:deney-donanim-kartlar} üzerindeki 1--14 numaraları doğrudan Tablo~\ref{tab:deney-alt-birimler} içindeki açıklamalara karşılık gelmektedir.

\begin{figure}[H]
  \centering
  \includegraphics[width=0.72\textwidth]{deneysel/onceki_tez_sekil_5_3_donanim_alt_birimleri.jpeg}
  \caption{Evirici donanımı ve alt birimlerinin numaralı tanıtımı.}
  \label{fig:deney-donanim-kartlar}
\end{figure}

\clearpage
\begin{table}[H]
  \centering
  \caption{Evirici donanımının alt birimleri.}
  \label{tab:deney-alt-birimler}
  \small
  \begin{tabularx}{0.94\textwidth}{X X}
    \toprule
    1. LCL filtrelerin $L_1$ bobinleri & 2. IGBT kapı sürme kartları \\
    3. Nötr IGBT'leri & 4. Yükseltici (boost) diyotları \\
    5. Fonksiyon değiştirme röleleri & 6. Snubber kapasitörleri \\
    7. Kontrol kartı & 8. Tampon (buffer) kartı \\
    9. Ölçme kartı & 10. LCL filtre kapasitörleri \\
    11. Ani akım (inrush) röleleri & 12. Kumanda beslemeleri \\
    13. Şarj fonksiyonu lambası & 14. Deşarj fonksiyonu lambası \\
    \bottomrule
  \end{tabularx}
\end{table}

Şekil~\ref{fig:deney-faz-donanim-1} ve Şekil~\ref{fig:deney-faz-donanim-2}, aynı faz kolunun iki farklı açıdan çekilmiş ayrıntılarını göstermektedir; her iki görseldeki 1--8 numaraları Tablo~\ref{tab:deney-faz-komponentler} ile birlikte okunmalıdır. İki görüntünün ayrı tutulması; boost diyotları, fonksiyon değiştirme röleleri, DA bara ve snubber kapasitörleri, soğutma düzeni, sürücü kartı ve $L_1$ bobininin fiziksel konumunu kaybetmeden izlemeyi sağlamaktadır.

\begin{figure}[H]
  \centering
  \includegraphics[height=0.68\textheight]{deneysel/onceki_tez_sekil_5_4_faz_donanimlari_1.png}
  \caption{Eviricinin bir fazına ait donanımlar, birinci görünüm.}
  \label{fig:deney-faz-donanim-1}
\end{figure}

\begin{figure}[H]
  \centering
  \includegraphics[height=0.68\textheight]{deneysel/onceki_tez_sekil_5_5_faz_donanimlari_2.png}
  \caption{Eviricinin bir fazına ait donanımlar, ikinci görünüm.}
  \label{fig:deney-faz-donanim-2}
\end{figure}

\clearpage
\begin{table}[H]
  \centering
  \caption{Şekil~\ref{fig:deney-faz-donanim-1} ve Şekil~\ref{fig:deney-faz-donanim-2}'deki faz donanımı bileşenleri.}
  \label{tab:deney-faz-komponentler}
  \small
  \begin{tabularx}{0.94\textwidth}{X X}
    \toprule
    1. Boost diyotları ve kondansatörü & 2. Fonksiyon değiştirme röleleri \\
    3. Snubber kondansatörleri & 4. DA bara elektrolitik kondansatörleri \\
    5. IGBT soğutucu & 6. Yüksek devirli soğutma fanı \\
    7. IGBT kapı sürme kartı & 8. LCL filtre $L_1$ bobini \\
    \bottomrule
  \end{tabularx}
\end{table}

Şekil~\ref{fig:deney-giris-devresi}, evirici giriş devresini ve 1--6 ile işaretlenmiş yardımcı donanımları göstermektedir; numaraların karşılığı Tablo~\ref{tab:deney-giris-komponentler} içindedir. Giriş bölümünün ayrı verilmesi; inrush röleleri, sigorta, yardımcı kaynaklar, EMI filtresi, kontrol kartı ve tampon kartının güç katıyla bağlantısını görünür kılmaktadır.

\begin{figure}[H]
  \centering
  \includegraphics[width=0.94\textwidth]{deneysel/onceki_tez_sekil_5_6_giris_devresi.png}
  \caption{Evirici giriş devresinin ayrıntı görünümü.}
  \label{fig:deney-giris-devresi}
\end{figure}

\clearpage
\begin{table}[H]
  \centering
  \caption{Evirici giriş kısmına ait bileşenler.}
  \label{tab:deney-giris-komponentler}
  \small
  \begin{tabularx}{0.94\textwidth}{X X}
    \toprule
    1. Ani akım (inrush) röleleri & 2. Evirici açma kapama sigortası \\
    3. 12 V--24 V güç kaynakları & 4. EMI filtre \\
    5. Kontrol kartı & 6. Tampon (buffer) kartı \\
    \bottomrule
  \end{tabularx}
\end{table}

\section{Kontrol Kartı, Tampon Kartı ve Kapı Sürücüsü}

Şekil~\ref{fig:deney-kontrol-buffer}, kontrol ve tampon kartlarını numaralı işlevleriyle göstermektedir. Kontrol kartında RS485 haberleşme birimi, optoizole giriş/çıkışlar, DA bara ve akım ölçüm dönüştürme devreleri, PIC32MK ailesi mikrodenetleyici, lineer regülatörler ve röle/fan çıkış katları bulunmaktadır. Tampon kartı, mikrodenetleyicinin \SI{3.3}{\volt} tek uçlu (single-ended) PWM işaretlerini diferansiyel sinyal çiftlerine dönüştürmek ve uzun ekranlı kablolarda ortak-mod gürültü etkisini azaltmak amacıyla kullanılmıştır. Fotoğraftaki numaralar sırasıyla aşağıdaki işlevleri gösterir: (1) RS485 haberleşme; (2)--(3) optoizole giriş ve çıkışlar; (4) diferansiyel DA bara gerilimini tek uçlu biçime dönüştüren ölçüm katı; (5) ölçüm, haberleşme ve kontrol algoritmalarını yürüten PIC32MK mikrodenetleyici; (6) akım ölçüm dönüştürme katı; (7) \SI{5}{\volt} lineer regülatörler; (8) röle ve fanlar için MOSFET güç çıkışları; (9) \SI{9}{\volt} anahtarlamalı regülatör; (10) PWM tampon kartının \SI{9}{\volt} beslemesi; (11) tek uçlu PWM işaretlerini diferansiyel çiftlere dönüştüren kat.

\begin{figure}[H]
  \centering
  \includegraphics[width=0.90\textwidth]{deneysel/onceki_tez_sekil_5_7_kontrol_buffer.png}
  \caption{Kontrol kartı ve tampon kartının numaralı görünümü.}
  \label{fig:deney-kontrol-buffer}
\end{figure}

PWM1--PWM6 kanallarının yüksek ve alçak işaretleri tampon kartında diferansiyel çiftlere dönüştürülmekte; yaklaşık \SI{12}{\volt} genlikte, tek yönlü iletim amacıyla H-köprü tabanlı bir sürme çözümü kullanılmaktadır. PWM zamanlayıcısındaki ölü zaman tasarım değeri \SI{1}{\micro\second}'dir (Tablo~\ref{tab:sistem-parametreleri}).

Şekil~\ref{fig:deney-igbt-surucu}, ADuM4135 tabanlı IGBT kapı sürme kartını göstermektedir. Diferansiyel giriş, desatürasyon koruması, aktif Miller kenetleme ve düşük gerilim kilitleme özellikleri sürücü seçiminin gerekçeleridir: diferansiyel giriş uzun PWM hatlarındaki ortak-mod gürültüsüne bağışıklık sağlar; desatürasyon koruması aşırı akımda IGBT komutunu donanımsal olarak keser; aktif Miller kenetleme yüksek frekanslı anahtarlamada hatalı iletime geçme (false turn-on) riskini azaltır. Bütün sürücülerden gelen FAULT çıkışları kontrol kartında toplanmakta ve ortak RESET işareti sürücülere paralel dağıtılmaktadır. Sonuç olarak donanım koruması ile denetim yazılımı arasında bir hata yönetim zinciri oluşmaktadır.

\begin{figure}[H]
  \centering
  \includegraphics[height=0.64\textheight]{deneysel/onceki_tez_sekil_5_8_igbt_surucu.jpeg}
  \caption{IGBT kapı sürme devresi.}
  \label{fig:deney-igbt-surucu}
\end{figure}

\clearpage
\section{Gerilim ve Akım Ölçme Devreleri}

Faz--nötr gerilimleri ile pozitif ve negatif DA bara gerilimleri, AMC1200 tabanlı izole ölçüm devreleriyle alınmaktadır; ölçüm kartının ofset, ADC aralığı ve DA bara ölçüm aralığı Tablo~\ref{tab:sistem-parametreleri} içindedir. AMC1200'ün gerilim kazancı 8 olduğundan girişte izin verilen en yüksek genlik \SI{162.5}{\milli\volt} olarak hesaplanmış; faz ve DA bara ölçümleri için ayrı gerilim bölücüler kullanılmıştır. Faz ölçüm sınırı anlık yaklaşık \SI{387}{\volt}'tur.

Mikrodenetleyicideki yedi ADC kanalının üçü faz--nötr gerilimlerine, ikisi pozitif ve negatif DA bara gerilimlerine ayrılmıştır. Faz ölçümleri AMC1200 çıkışlarından ekranlı kabloyla diferansiyel olarak taşınmış ve RC örtüşme önleme (anti-aliasing) filtrelerinden geçirilmiştir. DA bara kanallarında diferansiyel işaretler enstrümantasyon yükselteciyle işlenip \SI{1.3}{\volt} ofset eklenerek tek uçlu biçime dönüştürülmektedir. 12 bit ADC gösterimi, diferansiyel sayısal aralığı yaklaşık $-2047$ ile $+2048$ arasında tanımlamaktadır.

Dört telli T-tipi yapıda nötr noktasının kayması iki DA bara kondansatörünün gerilimlerini eşitsizleştirebileceğinden pozitif ve negatif bara kanalları sürekli izlenmekte; eşik aşımında mikrodenetleyici FAULT üretip PWM çıkışlarını kesmektedir. AMC1200 giriş beslemesi RFM-0505S izole dönüştürücü ile sağlanmış; dönüştürücü çıkışındaki regülasyon zayıflığını azaltmak için TCR2EF250 düşük gerilim düşümlü regülatör kullanılmıştır. Ölçüm zinciri, ölçeklemenin yanında izolasyon ve nötr nokta koruması görevlerini de üstlenmektedir.

Akım ölçümünde HLSR 20-P/SP33 Hall etkili sensör kullanılmıştır. Sensör kazancı ve ölçülebilir simetrik tepe akımı Tablo~\ref{tab:sistem-parametreleri} içindedir; tepe sınırın RMS karşılığı yaklaşık \SI{39.9}{\ampere}'dir. Faz başına akım tasarım değeri ve yüzde 20 paylı üst sınırı aynı tabloda verilmiştir. Hesaplar nominal direnç değerlerine dayanmaktadır.

Ani faz ve kondansatör akımlarının sınırlandırılması amacıyla kullanılan pencere karşılaştırıcı devre Şekil~\ref{fig:deney-pencere-karsilastirici} üzerinde verilmiştir. Devrede üst eşik \SI{2.0}{\volt}, alt eşik \SI{0.6}{\volt} olarak, sensör referansına simetrik ayarlanmıştır. Eşiklerin referanstan uzaklığı, sensör kazancıyla Tablo~\ref{tab:sistem-parametreleri} içindeki aşırı akım eşiğine karşılık gelmektedir. Eşik dışına çıkıldığında FAULT işareti kontrol kartında harici kesme oluşturarak PWM çıkışlarını kesmektedir.

\begin{figure}[H]
  \centering
  \includegraphics[width=0.96\textwidth]{deneysel/onceki_tez_sekil_5_9_pencere_karsilastirici.jpeg}
  \caption{Pencere karşılaştırıcı devresi.}
  \label{fig:deney-pencere-karsilastirici}
\end{figure}

Şekil~\ref{fig:deney-pencere-dalga}, harici bir deneme kartının DAC çıkışıyla uygulanan giriş işareti ve pencere karşılaştırıcının FAULT çıkışını göstermektedir. İmleç seviyeleri \SI{2.0}{\volt} ve \SI{0.6}{\volt} olarak ayarlanmış; giriş sınırların dışına çıktığında ikinci kanal işaretinin lojik sıfıra indiği gözlenmiştir. Kayıt, koruma mantığının işlevsel örneğidir; gecikme süresi ve gerçek kısa devre kesme enerjisi aynı görüntüden çıkarılamaz.

\begin{figure}[H]
  \centering
  \includegraphics[width=0.88\textwidth]{deneysel/onceki_tez_sekil_5_10_pencere_dalga.jpeg}
  \caption{Pencere karşılaştırıcı giriş ve çıkış gerilim dalga şekilleri.}
  \label{fig:deney-pencere-dalga}
\end{figure}

Şekil~\ref{fig:deney-olcum-karti}; HLSR akım sensörleri, AMC1200 izole ölçüm kanalları ve akım probu bağlantı iletkenlerini içeren ölçüm kartını göstermektedir. Kart üzerindeki \SI{4}{\milli\metre\squared} kesitli ters-U biçimli iletkenler; faz başına $L_1$ akımı, kondansatör akımı ve faz akımının ayrı ayrı osiloskop probuyla ölçülmesini sağlamaktadır. TLV431 referansları ve MC33269 regülatörleri ölçüm beslemesi ile \SI{1.3}{\volt} ofset üretiminde kullanılmıştır.

\begin{figure}[H]
  \centering
  \includegraphics[width=0.96\textwidth]{deneysel/onceki_tez_sekil_5_11_olcum_karti.jpeg}
  \caption{Ölçüm kartı.}
  \label{fig:deney-olcum-karti}
\end{figure}

\section{DA ve AA Güç Analizörleri}

DA bara gerilimi, akımı ve gücü Accuenergy AcuDC240 analizörüyle ölçülmüştür; üretici akım ve gerilim için yüzde 0,2, güç ve enerji için yüzde 0,5 hata toleransı bildirmektedir. Şekil~\ref{fig:deney-acudc240}(a) cihazın pano üzerindeki konumunu, Şekil~\ref{fig:deney-acudc240}(b) üretici dokümanındaki şönt bağlantısını göstermektedir.

\begin{figure}[H]
  \centering
  \begin{subfigure}[t]{0.46\textwidth}
    \centering
    \includegraphics[width=\linewidth]{deneysel/onceki_tez_sekil_5_12a_acudc240.jpeg}
    \caption{Evirici donanımı üzerindeki AcuDC240.}
  \end{subfigure}\hfill
  \begin{subfigure}[t]{0.50\textwidth}
    \centering
    \includegraphics[width=\linewidth]{deneysel/onceki_tez_sekil_5_12b_acudc240_baglanti.jpeg}
    \caption{DA analizörü ve şönt bağlantı şekli.}
  \end{subfigure}
  \caption{AcuDC240 DA bara analizörü ve bağlantısı.}
  \label{fig:deney-acudc240}
\end{figure}

Prototipte şönt direnci DA baranın negatif terminaline bağlanarak alçak-taraf akım ölçümü yapılmıştır. Şekil~\ref{fig:deney-sont-baglanti}, \SI{20}{\ampere}/\SI{60}{\milli\volt} şönt direncin fiziksel bağlantısını göstermektedir. Şöntün ölçüm uçları koaksiyel ve ekranlı kabloyla analizöre taşınmış; DA gerilim girişinde yüksek izolasyon dayanımlı kablo kullanılmıştır.

\begin{figure}[H]
  \centering
  \includegraphics[width=0.94\textwidth]{deneysel/onceki_tez_sekil_5_13_sont_baglanti.jpeg}
  \caption{\SI{20}{\ampere}--\SI{60}{\milli\volt} şönt direnç bağlantısı.}
  \label{fig:deney-sont-baglanti}
\end{figure}

AA tarafta Entes EMK-01 analizörü ve 50/5 oranlı, \SI{5}{\volt\ampere}, Class~0.1 akım transformatörleri kullanılmıştır. Akım transformatörünün IEC~61869-2'ye göre yapılan kalibrasyon ölçümü EK A'da verilmiştir; deney akımının bulunduğu anma akımının \%5'i bölgesinde oran hatası $-$\%0,046, faz hatası 2,8~dakikadır ve AA güç ölçümünün belirsizliğinde ihmal edilebilir düzeydedir. Üç faz ve nötr hattı için dört akım transformatörü bağlanmış; faz--nötr gerilimleri, faz/nötr akımları, aktif ve reaktif güç ile toplam THD değerleri Modbus RTU üzerinden okunmuştur. Şekil~\ref{fig:deney-emk01}(a) analizörü, Şekil~\ref{fig:deney-emk01}(b) üreticinin bağlantı düzenini vermektedir.

\begin{figure}[H]
  \centering
  \begin{subfigure}[t]{0.43\textwidth}
    \centering
    \includegraphics[width=\linewidth]{deneysel/onceki_tez_sekil_5_14a_emk01.jpeg}
    \caption{EMK-01 AA analizörü.}
  \end{subfigure}\hfill
  \begin{subfigure}[t]{0.53\textwidth}
    \centering
    \includegraphics[width=\linewidth]{deneysel/onceki_tez_sekil_5_14b_emk01_baglanti.png}
    \caption{Üretici tarafından önerilen bağlantı şekli.}
  \end{subfigure}
  \caption{EMK-01 AA analizörü ve bağlantısı.}
  \label{fig:deney-emk01}
\end{figure}

\section{PC Ana Kumanda Yazılımı ve Veri Toplama}

PC ana kumanda yazılımı; EMK-01, AcuDC240 ve kontrol kartından RS485/Modbus RTU üzerinden veri almak, anlık değerleri grafikleştirmek ve verim hesabına temel oluşturmak amacıyla geliştirilmiştir. Şekil~\ref{fig:deney-pc-arayuzu} yazılımın arayüzünü göstermektedir. Bir veri çevrimi yaklaşık \SI{1}{\second} sürmektedir. Her çevrimde önce DA analizörü, en fazla yaklaşık \SI{50}{\milli\second} sonra AA analizörü okunmaktadır. Kayıt yapısı kararlı durum eğilimlerini izlemek ve verimi hesaplamak için kullanılmıştır.

Haberleşme hızı cihazların izin verdiği en yüksek değer olan 38400 baud seçilmiş; veri sırası DA analizörü, AA analizörü ve kontrol kartı biçiminde kurulmuştur. Yazılımın karşılaştırma seti verim, THD, güç faktörü, faz akımı, faz--nötr gerilimi ve DA bara geriliminden oluşmaktadır. Osiloskop kanal adlarında ``kontrolör çıkışı'' ifadesi geçmektedir. Senaryo tanımlarında S1--S3 açık çevrim modülasyonla çalışır ve akım geri beslemesi yalnız S4'te bulunur (Alt Bölüm~\ref{sec:s4-kapali-cevrim}); kanal adındaki ifade, benzetim modelinin yapısına göre modülasyon komutu olarak okunmuştur.

\begin{figure}[H]
  \centering
  \includegraphics[width=0.98\textwidth]{deneysel/onceki_tez_sekil_5_15_pc_arayuzu.png}
  \caption{Haberleşme sistemi arayüzü.}
  \label{fig:deney-pc-arayuzu}
\end{figure}

Ortak hedef çalışma koşulu olarak DA bara gerilimi yaklaşık \SI{400}{\volt}, faz--nötr gerilimi yaklaşık \SI{130}{\volt}, faz akımı yaklaşık \SI{2.3}{\ampere} ve birim güç faktörüne yakın işletim tanımlanmıştır. Üç deneyde gerçekleşen DA bara gerilimleri (\num{431.5}, \num{395.0} ve \SI{427.5}{\volt}) ve şebeke gerilimleri (\num{119.8}, \num{121.0} ve \SI{124.3}{\volt}) birbirinden farklıdır (Tablo~\ref{tab:gelen-excel-senaryolar}). Her benzetim kendi deneyinin DA bara ve şebeke gerilimiyle yapılmıştır. Verim ve THD, Eşitlik~\eqref{eq:verim-tanimi} ve Eşitlik~\eqref{eq:thd-rmse-tanimlari} ile hesaplanmıştır.

\section{Klasik Anahtarlama Modu Deney Kayıtları}

Klasik yöntemde bütün anahtarlar tamamlayıcı sürülmekte ve düşük akımda senkron-buck benzeri ters akım yolu açık kalmaktadır. Deney özeti yüzde 92,5 verim, yüzde 15 faz akımı THD'si, yaklaşık \SI{2.3}{\ampere} faz akımı, \SI{120}{\volt} faz--nötr gerilimi ve \SI{430}{\volt} DA bara gerilimi bildirmektedir.

Şekil~\ref{fig:deney-klasik-pc}, PC yazılımında klasik anahtarlama sırasında elde edilen grafiklerin tamamını göstermektedir; AA/DA analizör değerleri ile güç ve THD eğilimleri aynı arayüzde izlenmektedir. THDi eğrisi yaklaşık 100 saniyelik kayıt boyunca \%5--22 arasında dalgalanmaktadır; özetteki \%15 kaydın ortalama düzeyini, ekrandaki \%11,1 ise tek bir anlık okumayı göstermektedir.

\begin{figure}[H]
  \centering
  \includegraphics[width=0.98\textwidth]{deneysel/onceki_tez_sekil_5_16_klasik_pc.jpeg}
  \caption{Klasik anahtarlama sırasında PC yazılımında alınan grafikler.}
  \label{fig:deney-klasik-pc}
\end{figure}

Şekil~\ref{fig:deney-klasik-l1} üzerinde yukarıdan aşağıya kontrol işareti, şebeke gerilimi, işlemciden alınan ortalama endüktör akımı ve akım probuyla ölçülen $L_1$ akımı bulunmaktadır. Kayıt, anahtarlama dalgalanmasını ve akımın sıfır çevresindeki davranışını doğrudan göstermektedir.

\begin{figure}[H]
  \centering
  \includegraphics[width=0.94\textwidth]{deneysel/onceki_tez_sekil_5_17_klasik_l1.png}
  \caption{Klasik yöntemde kontrol işareti, şebeke gerilimi, ortalama endüktör akımı ve $L_1$ akımı.}
  \label{fig:deney-klasik-l1}
\end{figure}

Şekil~\ref{fig:deney-klasik-faz} üzerinde ilk üç kanal korunmuş, dördüncü kanalda faz akımı verilmiştir. $L_1$ akımı ile faz akımının ayrı kayıtlarda sunulması, filtre ve şebeke tarafındaki dalga biçimlerinin aynı fiziksel büyüklük olarak yorumlanmasını önlemektedir. Faz akımının tepesi yuvarlaktır ve tepeyi bir kambur izlemektedir; izin sayısallaştırılmasıyla elde edilen THD yaklaşık \%21'dir (Alt Bölüm~\ref{sec:thd-kaynaklar}).

\begin{figure}[H]
  \centering
  \includegraphics[width=0.94\textwidth]{deneysel/onceki_tez_sekil_5_18_klasik_faz.png}
  \caption{Klasik yöntemde kontrol işareti, şebeke gerilimi, ortalama endüktör akımı ve faz akımı.}
  \label{fig:deney-klasik-faz}
\end{figure}

\section{Nötr Anahtar Susturma Modu Deney Kayıtları}

Nötr anahtar susturma yönteminde ilgili yarım periyotta ters akım yolu kesilmekte, düşük görev oranında kesintili iletim ve üçüncü harmonik ağırlıklı akım zarfı ortaya çıkmaktadır. Deney özeti yüzde 94,2 verim, yüzde 47 faz akımı THD'si, yaklaşık \SI{2.3}{\ampere} faz akımı, \SI{120}{\volt} faz--nötr gerilimi ve \SI{395}{\volt} DA bara gerilimi vermektedir.

\begin{figure}[H]
  \centering
  \includegraphics[width=0.98\textwidth]{deneysel/onceki_tez_sekil_5_19_susturma_pc.jpeg}
  \caption{Nötr anahtar susturma modunda PC yazılımında elde edilen grafikler.}
  \label{fig:deney-susturma-pc}
\end{figure}

Şekil~\ref{fig:deney-susturma-l1}, anahtarlamalı endüktör akımının sıfır çevresinde kesildiği aralıkları göstermektedir. Gözlenen zarf, Bölüm~\ref{chap:anahtarlama} içindeki kalibre S2 benzetimi ve üçüncü harmonik mekanizmasıyla nitel olarak uyumludur.

\begin{figure}[H]
  \centering
  \includegraphics[width=0.94\textwidth]{deneysel/onceki_tez_sekil_5_20_susturma_l1.png}
  \caption{Susturma yönteminde kontrol işareti, şebeke gerilimi, ortalama endüktör akımı ve $L_1$ akımı.}
  \label{fig:deney-susturma-l1}
\end{figure}

Şekil~\ref{fig:deney-susturma-faz}, aynı yöntemin faz akımındaki belirgin biçim bozulmasını göstermektedir.

\begin{figure}[H]
  \centering
  \includegraphics[width=0.94\textwidth]{deneysel/onceki_tez_sekil_5_21_susturma_faz.png}
  \caption{Susturma yönteminde kontrol işareti, şebeke gerilimi, ortalama endüktör akımı ve faz akımı.}
  \label{fig:deney-susturma-faz}
\end{figure}

\section{Susturma ve Sin\texorpdfstring{\textsuperscript{3}}{3} Modu Deney Kayıtları}

Üçüncü yöntemde nötr anahtar susturma korunmuş ve akım zarfındaki harmonik bozulmayı azaltmak amacıyla Sin\textsuperscript{3} tabanlı terim eklenmiştir. Deney özeti yüzde 94,5 verim, yüzde 14 faz akımı THD'si, yaklaşık \SI{2.3}{\ampere} faz akımı, \SI{124}{\volt} faz--nötr gerilimi ve \SI{427}{\volt} DA bara gerilimi bildirmektedir. Benzetim modelinde terim, deneydeki uygulamaya karşılık gelecek biçimde $K_{33}\sin^3\theta$ olarak tanımlanmıştır (Bölüm~\ref{chap:sin3}).

\begin{figure}[H]
  \centering
  \includegraphics[width=0.98\textwidth]{deneysel/onceki_tez_sekil_5_22_sin3_pc.jpeg}
  \caption{Sin\textsuperscript{3} uygulanmış susturma modunda PC yazılımında alınan grafikler.}
  \label{fig:deney-sin3-pc}
\end{figure}

Şekil~\ref{fig:deney-sin3-l1}, susturma deneyine göre daha düzenli bir $L_1$ akım zarfı göstermektedir; nitel iyileşme, bildirilen THD azalmasıyla uyumludur.

\begin{figure}[H]
  \centering
  \includegraphics[width=0.94\textwidth]{deneysel/onceki_tez_sekil_5_23_sin3_l1.png}
  \caption{Sin\textsuperscript{3} uygulanmış yöntemde kontrol işareti, şebeke gerilimi, ortalama endüktör akımı ve $L_1$ akımı.}
  \label{fig:deney-sin3-l1}
\end{figure}

Şekil~\ref{fig:deney-sin3-faz}, faz akımındaki zarfın susturma yöntemine göre düzeldiğini göstermektedir.

\begin{figure}[H]
  \centering
  \includegraphics[width=0.94\textwidth]{deneysel/onceki_tez_sekil_5_24_sin3_faz.png}
  \caption{Sin\textsuperscript{3} uygulanmış yöntemde kontrol işareti, şebeke gerilimi, ortalama endüktör akımı ve faz akımı.}
  \label{fig:deney-sin3-faz}
\end{figure}

\section{Deney Sonuçlarının Değerlendirilmesi}

Üç yöntemin deneysel ve kalibre benzetim verim ve faz akımı THD değerleri Tablo~\ref{tab:deney-dusuk-guc-ozeti} üzerindedir. AA analizörü şebeke tarafındaki faz akımını ölçtüğünden tabloda yalnız faz akımı THD'si yer almaktadır. Kalibrasyon yöntemi ve bağımsız doğrulama ölçütleri Alt Bölüm~\ref{sec:kalibrasyon} altında ayrıca verilmektedir.

\begin{table}[H]
  \centering
  \caption{Üç yöntemin deney özeti ve kalibre benzetimdeki verim ve faz akımı THD değerleri.}
  \label{tab:deney-dusuk-guc-ozeti}
  \footnotesize
  \begin{tabularx}{\textwidth}{X *{4}{>{\centering\arraybackslash}p{0.13\textwidth}}}
    \toprule
    & \multicolumn{2}{c}{Deney} & \multicolumn{2}{c}{Kalibre benzetim} \\
    \cmidrule(lr){2-3}\cmidrule(lr){4-5}
    Yöntem & Verim & THD$_{i,\mathrm{faz}}$ & Verim & THD$_{i,\mathrm{faz}}$ \\
    \midrule
    Klasik tamamlayıcı anahtarlama & \%92,50 & \%15 & \%91,15 & \%39,5 \\
    Nötr anahtar susturma & \%94,20 & \%47 & \%93,93 & \%56,7 \\
    Susturma + Sin\textsuperscript{3} & \%94,50 & \%14 & \%94,25 & \%14,8 \\
    \bottomrule
  \end{tabularx}
\end{table}

Deneyde susturma verimi 1,7 puan artırırken THD'yi \%47'ye çıkarmıştır. Sin\textsuperscript{3} terimi ise verimi \%94,5 düzeyinde korurken THD'yi \%47'den \%14'e, klasik anahtarlama düzeyine indirmiştir. Kalibre benzetim, verim ve THD sıralamasını deneyle aynı vermektedir.

Deney düzeneğindeki güç, ölçüm ve veri toplama akışı Şekil~\ref{fig:deney-olcum-zinciri} üzerinde özetlenmiştir.

\begin{figure}[H]
  \centering
  \begin{tikzpicture}[
    node distance=8mm and 10mm,
    block/.style={draw=blue!55!black,fill=blue!7,rounded corners,align=center,minimum height=10mm,text width=28mm,font=\small},
    meas/.style={draw=green!50!black,fill=green!8,rounded corners,align=center,minimum height=10mm,text width=28mm,font=\small},
    data/.style={draw=orange!70!black,fill=orange!10,rounded corners,align=center,minimum height=10mm,text width=30mm,font=\small},
    arr/.style={-{Latex[length=2.5mm]},thick}
  ]
    \node[block] (dc) {DA bara\\AcuDC240 + şönt};
    \node[block,right=of dc] (inv) {T-tipi evirici\\LCL filtre};
    \node[block,right=of inv] (grid) {Varyak ve\\AA şebeke};
    \draw[arr] (dc) -- node[above,font=\scriptsize] {$P_{\mathrm{dc}}$} (inv);
    \draw[arr] (inv) -- node[above,font=\scriptsize] {$v_g,i_g$} (grid);
    \node[meas,below=of inv] (scope) {Osiloskop\\$i_{L1}$, $i_g$, $v_g$, DAC};
    \node[meas,below=of grid] (ac) {EMK-01 + CT\\$P_{\mathrm{ac}}$, THD, PF};
    \node[meas,below=of dc] (ctrl) {Kontrol/ölçüm kartı\\ADC, PWM, FAULT};
    \draw[arr] (inv) -- (scope);
    \draw[arr] (grid) -- (ac);
    \draw[arr] (inv) -- (ctrl);
    \node[data,below=15mm of scope] (pc) {PC ana kumanda yazılımı\\RS485/Modbus RTU};
    \draw[arr] (ctrl) -- (pc);
    \draw[arr] (ac) -- (pc);
    \draw[arr] (dc) |- (pc);
    \draw[arr,dashed] (scope) -- node[right,font=\scriptsize] {ekran kaydı} (pc);
  \end{tikzpicture}
  \caption{Deney düzeneğindeki güç, ölçüm ve veri toplama akışı.}
  \label{fig:deney-olcum-zinciri}
\end{figure}

Tablo~\ref{tab:olcum-cihazlari}, kullanılan cihaz ve sensörleri bildirilen özellikleri ve doğruluk bilgisinin kaynağıyla birlikte toplamaktadır. AcuDC240 için hesaplarda veri sayfasındaki yüzde hata değerleri kullanılmıştır.

Üç yöntemin çalışma koşulları Tablo~\ref{tab:deney-kosullari} üzerinde karşılaştırılmıştır.

Deney özet değerleri ile PC ekranlarında görünen anlık değerler arasındaki farklar Tablo~\ref{tab:deney-metin-ekran-karsilastirma} üzerinde verilmiştir. Ekran değeri görselde okunabilen tek anlık değer, özet değeri ise deney raporundaki sayıdır. Özet THDi değeri süreli kaydın ortalama düzeyidir. Tablo, kayıtların kendi içindeki tutarlılığını göstermektedir.

Deneyler yaklaşık \SI{2.3}{\ampere} faz akımındaki düşük güç noktasında yapılmıştır. Orta ve anma akımında aynı denetim yapısının davranışı S4 benzetim taramasıyla (Alt Bölüm~\ref{sec:s4-kapali-cevrim}) incelenmiştir. Tablo~\ref{tab:deney-juri-durum}, her güç bölgesinin hangi yöntemle incelendiğini ve elde edilen sonucu özetlemektedir.

Anma gücü, donanımın tasarım hedefidir. S4 benzetiminde \SI{25}{\ampere} için büyütülen DA bara kapasitansı faz başına modelden kaynaklanır; üç fazlı donanımda toplam bara kapasitansı faz payının üç katıdır (Alt Bölüm~\ref{sec:s4-kapali-cevrim}). Deneysel verimin belirsizliği, DA analizörünün bildirilen yüzde 0,5 güç toleransından kestirilmiştir. Birleşik belirsizliğin $u_{\eta}=\eta\sqrt{(u_{P_{\mathrm{ac}}}/P_{\mathrm{ac}})^2+(u_{P_{\mathrm{dc}}}/P_{\mathrm{dc}})^2}$ ilk yaklaşımında DA tarafı katkısı \%94 verimde yaklaşık $\pm0{,}5$ puandır.

\begin{table}[H]
  \centering
  \caption{Ölçüm cihazları, bildirilen özellikleri ve doğruluk bilgisinin kaynağı.}
  \label{tab:olcum-cihazlari}
  \footnotesize
  \begin{tabularx}{\textwidth}{>{\raggedright\arraybackslash}p{0.16\textwidth} >{\raggedright\arraybackslash}p{0.21\textwidth} >{\raggedright\arraybackslash}p{0.20\textwidth} >{\raggedright\arraybackslash}X}
    \toprule
    Büyüklük & Cihaz / sensör & Bildirilen özellik & Doğruluk bilgisinin kaynağı \\
    \midrule
    AA güç ve THD & Entes EMK-01 ve 50/5 oranlı, \SI{5}{\volt\ampere} akım transformatörleri & Class 0.1 & Akım transformatörünün kalibrasyon ölçümü EK A'dadır (15.01.2025): anma akımının \%5'inde oran hatası $-$\%0,046, faz hatası 2,8~dk. \\
    DA güç & Accuenergy AcuDC240 ve \SI{20}{\ampere}/\SI{60}{\milli\volt} şönt & Akım/gerilim yüzde 0,2; güç/enerji yüzde 0,5 & Doğruluk değerleri veri sayfasıyla (EK A) doğrulanmıştır. \\
    Hızlı akım & Tektronix TCP0150 & \SI{150}{\ampere}, \SI{20}{\mega\hertz} & Bildirilen özellik üretici değeriyle aynıdır. \\
    Hızlı gerilim & Tektronix THDP0200 & \SI{1500}{\volt}, \SI{200}{\mega\hertz} & Bildirilen özellik üretici değeriyle aynıdır. \\
    Osiloskop & Tektronix DPO3034 & \SI{300}{\mega\hertz}, 2,5~GS/s & Bant genişliği ve örnekleme hızı üretici değeriyle aynıdır. \\
    Faz/LCL akımı & HLSR 20-P/SP33 & \SI{23}{\milli\volt\per\ampere}; \SI{39.9}{\ampere} RMS hesap sınırı & Duyarlılık veri sayfasından alınmıştır (EK A). \\
    Gerilim ölçümü & AMC1200 tabanlı izole kart & Kazanç 8; \SI{2.6}{\volt} ADC aralığı & Kazanç veri sayfasından alınmıştır (EK A). \\
    \bottomrule
  \end{tabularx}
\end{table}

\begin{table}[H]
  \centering
  \caption{Üç yöntem için bildirilen çalışma koşulları.}
  \label{tab:deney-kosullari}
  \footnotesize
  \begin{tabularx}{\textwidth}{p{0.21\textwidth} *{3}{>{\centering\arraybackslash}X} X}
    \toprule
    Büyüklük & Klasik & Nötr susturma & Susturma+Sin\textsuperscript{3} & Yorum \\
    \midrule
    Faz akımı & \SI{2.3}{\ampere} & \SI{2.3}{\ampere} & \SI{2.3}{\ampere} & Düşük güç noktası. \\
    Faz--nötr gerilimi & \SI{120}{\volt} & \SI{120}{\volt} & \SI{124}{\volt} & Deneyler arasında gerilim farkı vardır. \\
    DA bara gerilimi & \SI{430}{\volt} & \SI{395}{\volt} & \SI{427}{\volt} & Deneyler arasında bara gerilimi farkı vardır. \\
    Analizörün $\cos\varphi$ göstergesi & $-0{,}990$--$1{,}000$ & $-0{,}995$ & $-0{,}997$ & Temel bileşen güç faktörüdür ($\cos\varphi_1$); işaret yorumu Alt Bölüm~\ref{sec:kalibrasyon} içindedir. \\
    Analizörün THD$_v$ göstergesi & \%3,7 & \%5,4 & \%4,8 & Ölçüm varyak ile evirici arasında alınmıştır. \\
    \bottomrule
  \end{tabularx}
\end{table}

\begin{table}[H]
  \centering
  \caption{Deney özeti ile PC ekranlarında okunabilen anlık değerlerin karşılaştırılması.}
  \label{tab:deney-metin-ekran-karsilastirma}
  \footnotesize
  \begin{tabularx}{\textwidth}{p{0.21\textwidth} p{0.16\textwidth} >{\centering\arraybackslash}X >{\centering\arraybackslash}X X}
    \toprule
    Yöntem & Büyüklük & Özet & PC ekranı & Farkın durumu \\
    \midrule
    Klasik & Verim & \%92,5 & \%92,9 & Ekran değeri anlık okumadır. \\
    Klasik & THD$_i$ & \%15 & \%11,1 & Anlık okuma; eğri \%5--22 arasında dalgalanır. \\
    Klasik & $V_{\mathrm{dc}}$ & \SI{430}{\volt} & \SI{431.5}{\volt} & Yakın değerler. \\
    \midrule
    Nötr susturma & Verim & \%94,2 & \%93,8 & Ekran değeri anlık okumadır. \\
    Nötr susturma & THD$_i$ & \%47 & \%48,6 & Anlık okuma; eğri \%45--50 arasındadır. \\
    Nötr susturma & $V_{\mathrm{dc}}$ & \SI{395}{\volt} & \SI{394.7}{\volt} & Yakın değerler. \\
    \midrule
    Susturma+Sin\textsuperscript{3} & Verim & \%94,5 & \%93,9 & Ekran değeri anlık okumadır. \\
    Susturma+Sin\textsuperscript{3} & THD$_i$ & \%14 & \%15,7 & Anlık okuma; eğri kaydın büyük bölümünde \%14--18 arasındadır. \\
    Susturma+Sin\textsuperscript{3} & $V_{\mathrm{dc}}$ & \SI{427}{\volt} & \SI{427.5}{\volt} & Yakın değerler. \\
    \bottomrule
  \end{tabularx}
\end{table}

\begin{table}[H]
  \centering
  \caption{Güç bölgelerine göre inceleme yöntemi ve elde edilen sonuçlar.}
  \label{tab:deney-juri-durum}
  \small
  \begin{tabularx}{\textwidth}{p{0.15\textwidth} p{0.15\textwidth} X p{0.24\textwidth}}
    \toprule
    Çalışma bölgesi & İnceleme yöntemi & Kullanılan kayıt ve model & Elde edilen sonuç \\
    \midrule
    Düşük güç (\SI{2.3}{\ampere}) & Deney ve kalibre benzetim & Üç yönteme ait PC ve osiloskop ekranları ve deney özeti, S1--S3 kalibre benzetimi (Tablo~\ref{tab:kalibrasyon-karsilastirma}), S4 \SI{2.3}{\ampere} benzetimleri & Verim ve THD sıralaması deneyle aynıdır. S3 THD'si deneyde \%14, benzetimde \%14,8'dir. \\
    Orta güç (\SI{10}{\ampere}) & S4 benzetimi & S4 P, PI, PIR benzetimleri (Tablo~\ref{tab:s4-tarama}) & PIR ile temel akım genlik hatası $-$\%2,2, verim \%97,28. \\
    Anma (\SI{25}{\ampere}) & S4 benzetimi & S4 P, PI, PIR benzetimleri, büyütülmüş DA bara (Tablo~\ref{tab:s4-tarama}) & PIR ile temel akım genlik hatası $-$\%1,0, verim \%96,03. \\
    Geçici rejim & S4 benzetimi & S4 \SI{10}{\ampere}$\rightarrow$\SI{25}{\ampere} basamağı (Tablo~\ref{tab:s4-basamak}) & Aşım \%0,2, yerleşme süresi \SI{0.44}{\milli\second}. \\
    \bottomrule
  \end{tabularx}
\end{table}

\subsection{Deney Özet Tablosunun Sayısal Denetimi}

Deney özet tablosunun kaynağı olan çalışma kitabı, Senaryo~1--3 için birer satırlık özet değerler içermektedir. Tablo~\ref{tab:gelen-excel-senaryolar}, hücrelerdeki değerleri değiştirmeden ve $\eta_{VI}=100P_{\mathrm{out}}/(V_{\mathrm{dc}}I_{\mathrm{dc}})$ bağımsız hesabını ayrıca göstererek kaynağın izini korumaktadır.

\begin{table}[H]
  \centering
  \caption{Deney özet tablosundaki değerler ve DA giriş çarpımından yeniden hesaplanan verim.}
  \label{tab:gelen-excel-senaryolar}
  \footnotesize
  \begin{tabularx}{\textwidth}{c r r r r r r r}
    \toprule
    Senaryo & $V_{\mathrm{dc}}$ (V) & $I_{\mathrm{dc}}$ (A) & $V_{\mathrm{ac}}$ (V RMS) & $I_{faz}$ (A) & $P_o$ (W) & THD (\%) & $\eta_{VI}$ (\%) \\
    \midrule
    1 & 431,5 & 0,624 & 119,8 & 2,26 & 249,6 & 15 & 92,700 \\
    2 & 395,0 & 0,684 & 121,0 & 2,33 & 253,2 & 47 & 93,715 \\
    3 & 427,5 & 0,697 & 124,3 & 2,31 & 281,2 & 14 & 94,373 \\
    \bottomrule
  \end{tabularx}
\end{table}

Bildirilen verimler sırasıyla \%92,5, \%94,2 ve \%94,5'tir; yeniden hesaplanan değerlerle farklar yaklaşık $-0{,}20$, $+0{,}48$ ve $+0{,}13$ yüzde puandır. Tablo deneysel özet olarak kullanılmış, benzetim karşılaştırması Alt Bölüm~\ref{sec:kalibrasyon} altında modelin kendi benzetimlerinden yapılmıştır.

\subsection{Faz Akımı THD Değerlerinin Karşılaştırılması}
\label{sec:thd-kaynaklar}

Faz akımı THD'si için üç kaynak bulunmaktadır: AA analizörünün kaydı, osiloskop kayıtları ve kalibre benzetim. Osiloskop kayıtlarında (Şekil~\ref{fig:deney-klasik-faz}, Şekil~\ref{fig:deney-susturma-faz} ve Şekil~\ref{fig:deney-sin3-faz}) dördüncü kanaldaki faz akımı izi sayısallaştırılmış, iki ila dört tam şebeke periyodu üzerinde 25. harmoniğe kadar Fourier analizi yapılmıştır. Tablo~\ref{tab:thd-kaynaklar} sonuçları özetlemektedir. Ekran görüntüsünden yapılan hesabın düşey çözünürlüğü yaklaşık \SI{0.1}{\ampere}'dir. Osiloskop kayıtlarının çevrim RMS akımı S1'de \SI{2.87}{\ampere}, S2'de \SI{2.42}{\ampere}, S3'te \SI{2.29}{\ampere}'dir ve kayıtlar analizör özetinden ayrı anlarda alınmıştır.

\begin{table}[H]
  \centering
  \caption{Faz akımı THD değerlerinin kaynaklara göre karşılaştırılması.}
  \label{tab:thd-kaynaklar}
  \footnotesize
  \begin{tabularx}{\textwidth}{X *{4}{>{\centering\arraybackslash}p{0.15\textwidth}}}
    \toprule
    Yöntem & Analizör özeti & Analizör anlık okuması & Osiloskop izi & Kalibre benzetim \\
    \midrule
    Klasik & \%15 & \%11,1 & yaklaşık \%21 & \%39,5 \\
    Nötr susturma & \%47 & \%48,6 & yaklaşık \%47 & \%56,7 \\
    Susturma+Sin\textsuperscript{3} & \%14 & \%15,7 & yaklaşık \%12 & \%14,8 \\
    \bottomrule
  \end{tabularx}
\end{table}

Analizör ve osiloskop kayıtları üç yöntemde aynı büyüklük düzeyini vermektedir: S2 \%47 düzeyinde, S1 \%11--21, S3 \%12--16 aralığındadır. Osiloskop izinde üçüncü harmonik oranı S2'de yaklaşık \%47 iken S3'te yaklaşık \%9'dur; Sin\textsuperscript{3} teriminin hedeflediği bileşenin azalması deney kaydında doğrudan görülmektedir. Deneysel THD değeri olarak analizör özetindeki \%15, \%47 ve \%14 kullanılmıştır. Analizörün güç faktörü göstergesi tek bir anlık okumadır ve S1'de faz akımı kayıt boyunca yaklaşık \SI{2.0}{\ampere} ile \SI{2.3}{\ampere} arasında değişmektedir. Dolayısıyla THD değerlendirmesi analizör özetindeki THD değerlerine dayanmaktadır.

\section{Benzetim Modelinin Deney Özetine Kalibrasyonu}
\label{sec:kalibrasyon}

Deney özeti ile benzetim modelinin karşılaştırılabilmesi için S1--S3 benzetim modelleri deney özetinin çalışma noktalarına kalibre edilmiştir. Kalibrasyonun amacı, modelin deneyin hangi büyüklüklerini hangi hatayla yeniden ürettiğini göstermek ve kalibre benzetimi deney bulgusundan ayrı tutmaktır. Taban modeller, Alt Bölüm~\ref{sec:model-dogrulama-duzeltme} altında doğrulanan S1, S2 ve S3 modelleridir. Kalibre değerler Tablo~\ref{tab:sistem-parametreleri} içinde ``Kalibre'' türüyle verilmiştir. S4 ileri beslemesi kalibre S3 değerlerini kullanmaktadır.

\subsection{Yöntem}

Devre, yarı iletken ve filtre parametreleri kalibrasyonda değiştirilmemiştir. Her senaryoda yalnız iki serbest değişken vardır: modülasyon genlik ölçeği ($K_{1x}$; S3'te $K_{33}/K_{13}$ oranı sabit tutularak $K_{13}$ ile $K_{33}$ birlikte) ve şebekeye göre referans fazı $\varphi_x$. İki serbest değişken iki hedefe ayarlanmıştır: aktif güç $P_{\mathrm{ac}}$ deney değerinin $\pm$\%1'i içinde, temel bileşen reaktif gücünün büyüklüğü $|Q_{g,1}|$ ise deneydeki $\cos\varphi$ okumasının yuvarlama bandı içinde tutulmuştur. AA çıkış akımı $I_{\mathrm{ac}}$, verim, DA güç ve akımı, $|\cos\varphi_1|$ ve faz akımı THD'si kalibre edilmemiştir ve bağımsız doğrulama ölçütüdür. THD hedefi olarak deney özetinin \%15, \%47 ve \%14 değerleri alınmış, THD birincil ölçütler olan $P_{\mathrm{ac}}$, $I_{\mathrm{ac}}$ ve verimin ardından ikincil ölçüt olarak değerlendirilmiştir.

\textbf{Reaktif güç işareti ve ölçüm noktası.} Benzetimde güç büyüklükleri evirici çıkış uçlarında, yani EMC süzgecinden sonra ve $L_2$ ile temsil edilen şebeke ve varyak eşdeğerinden önce ölçülmektedir. Ölçüm noktası, deneyde AA analizörünün bağlı olduğu yaklaşık \SI{120}{\volt} tarafına karşılık gelmektedir. İşaret kuralı üreticidir: evirici$\rightarrow$şebeke yönü pozitiftir, $Q_{g,1}>0$ akımın gerilimin gerisinde olduğu ve eviricinin şebekeye reaktif güç verdiği durumdur. Benzetimdeki pozitif $Q_{g,1}$'in kaynağı filtre kapasitörüdür. Deney özetindeki reaktif güç $Q_o=P_o\tan(\arccos\cos\varphi)$ ile hesaplanmıştır ve eksi işareti analizörün $\cos\varphi$ göstergesinden gelir. Karşılaştırma, işaretten bağımsız büyüklükler ($|Q_{g,1}|$, $|\cos\varphi_1|$) üzerinden yapılmıştır. Deneydeki $\cos\varphi$ okuması, temel bileşen güç faktörüdür ($\cos\varphi_1$). Üç basamaklı $\cos\varphi$ okuması $|Q_{g,1}|$ için S1'de 0--7,9~var, S2'de 24,1--26,7~var, S3'te 19,9--23,6~var bandı verir. Şebeke gerilimi okumasının \SI{0.1}{\volt} çözünürlüğü $|Q_{g,1}|$'de ayrıca yaklaşık $\pm12$~var belirsizlik doğurur. Her senaryoda kalibrasyon öncesi modelin işaret dalı ($Q_{g,1}>0$) korunmuştur.

\textbf{Köprünün yumuşak kaynak davranışı.} İlk kalibrasyon adımları, evirici köprüsünün temel bileşen geriliminin modülasyonla bire bir kaymadığını göstermiştir. Köprü, iç gerilimi $V_{\mathrm{int}}=aV_0\angle\varphi$ ve iç empedansı $Z_{\mathrm{th}}$ olan bir Thevenin eşdeğeriyle temsil edilmiştir. Burada $a$ modülasyon genlik ölçeği, $V_0$ ölçek bire eşitken iç gerilimin genliği, $\varphi$ referans fazıdır ve köprü çıkış gerilimi $V_A=V_{\mathrm{int}}-Z_{\mathrm{th}}I_{L1}$ olmaktadır. Kalibrasyonun ardışık iki benzetiminden ölçülen iç empedans ve güç duyarlılıkları, kalibrasyon öncesi benzetimlerdeki ara seviye zaman oranıyla birlikte Tablo~\ref{tab:kopru-ic-empedansi} içinde verilmiştir.

\begin{table}[H]
  \centering
  \caption{Köprünün kalibrasyon benzetimlerinden ölçülen Thevenin iç empedansı, ara seviye zaman oranı (kalibrasyon öncesi benzetimler) ve aktif ile reaktif gücün serbest değişkenlere duyarlılığı.}
  \label{tab:kopru-ic-empedansi}
  \small
  \setlength{\tabcolsep}{4pt}
  \begin{tabular}{l c c c c c}
    \toprule
    Senaryo & $Z_{\mathrm{th}}$ (\si{\ohm}) & $|Z_{\mathrm{th}}|$ (\si{\ohm}) & Ara seviye (\%) & $\mathrm{d}P/\mathrm{d}\ln a$ (\si{\watt}) & $\mathrm{d}Q/\mathrm{d}\varphi$ (var/\si{\degree}) \\
    \midrule
    S1 & $1{,}18-j1{,}18$ & \num{1.67} & \num{0.5} & $\approx\num{10000}$ & $\approx-260$ \\
    S2 & $23{,}3+j16{,}7$ & \num{28.7} & \num{12.4} & $\approx\num{660}$ & $\approx-30$ \\
    S3 & $13{,}5-j2{,}2$ & \num{13.7} & \num{22.6} & $\approx\num{1400}$ & $-12\ldots-25$ \\
    \bottomrule
  \end{tabular}
\end{table}

S1'deki iç empedans, tamamlayıcı anahtarlamada akım dalgalanmasının her periyotta sıfırı kesmesiyle \SI{1}{\micro\second} ölü zamanın akıma bağlı gerilim hatası üretmesinden kaynaklanmaktadır. Kaba ölü zaman kestirimi $V_{\mathrm{dc}}t_df_s/\Delta i_L\approx\SI{431}{\volt}\cdot\SI{1}{\micro\second}\cdot\SI{10}{\kilo\hertz}/\SI{2.6}{\ampere}\approx\SI{1.7}{\ohm}$ ölçülen değerle aynı düzeydedir. S2 ve S3'te iç empedans bir mertebe büyüktür ve kaynağı süreksiz iletimdir. Susturmalı köprü düşük akımda akım kaynağına yaklaşmakta ve senkron olmayan buck eşdeğeriyle aynı davranışı göstermektedir (Alt Bölüm~\ref{sec:dusuk-akim-modelleme}). Thevenin empedansı LCL ağının \SI{50}{\hertz}'teki yaklaşık \SI{0.5}{\ohm} empedansından büyük olduğundan aktif güç esas olarak genlikle, reaktif güç esas olarak fazla değişmektedir. Genlik ile aktif güç ve faz ile reaktif güç arasındaki eşleme, sert gerilim kaynağı varsayımına dayanan ağ hesabının öngördüğü eşlemenin tersidir. İki değişken birbirini etkilediği için hedefler, benzetimlerden öğrenilen ampirik Jacobian ile iki boyutlu Newton yinelemesiyle birlikte çözülmüştür. S1 ve S2'de dört, S3'te altı kalibrasyon benzetimi gerekmiştir.

\subsection{Kalibrasyon sonucu ve bağımsız karşılaştırma}

Kalibrasyon S1'de genliği değiştirmemiş (\%0,0003), fazı $+0{,}074^{\circ}$ kaydırmıştır; parametre dosyasındaki genlik deneyin aktif gücünü \%0,96 içinde zaten vermektedir ve faz düzeltmesi yalnız $|Q_{g,1}|$'i bant içine taşır. S2'de genlik \%3,3, S3'te \%1,5 azalmış, fazlar $+0{,}11^{\circ}$ ve $+0{,}66^{\circ}$ kaymıştır. Genlik belirsizliği deney toleransının ölçülen duyarlılıklara bölünmesiyle S1'de $\pm$\%0,025, S2'de $\pm$\%0,4, S3'te $\pm$\%0,2; faz belirsizliği sırasıyla $\pm0{,}05^{\circ}$, $\pm0{,}4^{\circ}$ ve $\pm0{,}5$--$1^{\circ}$ bulunmuştur. S2 ve S3 fazları faz belirsizliğinin içindedir ve fiziksel bir gecikme olarak yorumlanmamıştır. S2 ve S3'teki genlik azalması, yaklaşık \SI{50}{\micro\second}'lik ortalama iletim süresinde darbe başına 0,7--\SI{1.6}{\micro\second} kısalmaya eşdeğerdir. Değer, sürücü ve anahtar gecikmelerinin darbe genişliğinden yediği süre mertebesindedir.

\begin{table}[H]
  \centering
  \caption{Deney özeti ile kalibre benzetim modelinin karşılaştırılması. İlk iki satır kalibrasyon hedefidir, öteki satırlar bağımsız karşılaştırmadır. Fark: \% göreli, ``p'' yüzde puan.}
  \label{tab:kalibrasyon-karsilastirma}
  \scriptsize
  \setlength{\tabcolsep}{2.6pt}
  \begin{tabular}{@{}l ccc ccc ccc@{}}
    \toprule
    & \multicolumn{3}{c}{S1 Klasik} & \multicolumn{3}{c}{S2 Susturma} & \multicolumn{3}{c}{S3 Susturma + Sin\textsuperscript{3}} \\
    \cmidrule(lr){2-4}\cmidrule(lr){5-7}\cmidrule(lr){8-10}
    Büyüklük & Deney & Benzetim & Fark & Deney & Benzetim & Fark & Deney & Benzetim & Fark \\
    \midrule
    $P_{\mathrm{ac}}$ (W) & 249,6 & 247,21 & $-$\%0,96 & 253,2 & 253,80 & $+$\%0,24 & 281,2 & 279,90 & $-$\%0,46 \\
    $|Q_{g,1}|$ (var) & 0--7,9 & 1,93 & bantta & 24,1--26,7 & 24,43 & bantta & 19,9--23,6 & 21,41 & bantta \\
    \midrule
    $I_{\mathrm{ac}}$ (A) & 2,26 & 2,218 & $-$\%1,8 & 2,33 & 2,416 & $+$\%3,7 & 2,31 & 2,279 & $-$\%1,4 \\
    $\eta$ (\%) & 92,5 & 91,15 & $-$1,35 p & 94,2 & 93,93 & $-$0,27 p & 94,5 & 94,25 & $-$0,25 p \\
    $V_{\mathrm{dc}}$ (V) & 431,5 & 431,48 & $\approx0$ & 395,0 & 394,95 & $\approx0$ & 427,5 & 427,45 & $\approx0$ \\
    $P_{\mathrm{dc}}$ (W) & 269,26 & 271,21 & $+$\%0,7 & 270,18 & 270,19 & $+$\%0,004 & 297,97 & 296,97 & $-$\%0,34 \\
    $P_{\mathrm{dc}}-P_{\mathrm{ac}}$ (W) & 19,66 & 24,00 & $+$\%22 & 16,98 & 16,39 & $-$\%3,5 & 16,77 & 17,07 & $+$\%1,8 \\
    $I_{\mathrm{dc}}$ (A) & 0,624 & 0,6286 & $+$\%0,7 & 0,684 & 0,6841 & $+$\%0,02 & 0,697 & 0,6947 & $-$\%0,32 \\
    THD$_{i,\mathrm{faz}}$ (\%) & 15 & 39,5 & $+$24,5 p & 47 & 56,7 & $+$9,7 p & 14 & 14,8 & $+$0,8 p \\
    $|\cos\varphi_1|$ & 1,000 & 0,99997 & $\approx0$ & 0,995 & 0,99538 & $+$0,0004 & 0,997 & 0,99708 & $+$0,0001 \\
    \bottomrule
  \end{tabular}
\end{table}

Tablo~\ref{tab:kalibrasyon-karsilastirma}, kalibre benzetimin deney özetiyle karşılaştırmasını vermektedir. İlk iki satır, iki serbest değişkenle ayarlanan aktif güç ve reaktif güç büyüklüğünü içeren kalibrasyon hedefidir. Öteki satırlar bağımsız karşılaştırmadır.

\textbf{S3.} Kalibre model deneyin AA akımını $-$\%1,4, verimini $-0{,}25$ puan, akım THD'sini $+0{,}8$ puan ve DA gücünü $-$\%0,3 hatayla yeniden üretmektedir. THD kalibre edilmemiştir: kalibrasyon öncesinde \%10,9 olan değer kalibrasyon sonrasında \%14,8 olmuştur (deney \%14). Sonuç, susturmalı çalışmanın düşük akımdaki dalga biçimi bozulmasının modelde doğru büyüklükte olduğunu gösteren en güçlü bağımsız kanıttır.

\textbf{S2.} Kalibre model DA gücünü ve DA akımını \%0,02 içinde, verimi $-0{,}27$ puan hatayla yeniden üretmektedir. Dalga biçimi deneyle aynı karakterdedir: alternans ortasında çöken ve sıfır geçişlerinin iki yanında tepe yapan zarf hem osiloskop kaydında (Şekil~\ref{fig:deney-susturma-faz}) hem modelde (Şekil~\ref{fig:s1-s2-kalibre-dalga}) görülmektedir. AA akımı $+$\%3,7, THD $+9{,}7$ puan farkla yeniden üretilmektedir. İki serbest değişken aktif ve reaktif güce ayrıldığından AA akımı ve THD bağımsız karşılaştırma olarak kalmaktadır. Farkın olası kaynakları, sıfır geçişi çevresindeki süreksiz iletim ayrıntısını belirleyen IGBT çıkış kapasitesi ile sürücü gecikmeleri, diyot toparlanma parametreleri ve şebeke gerilimi bozulmasıdır.

\textbf{S1.} Kalibre model AA akımını $-$\%1,8, verimi $-1{,}35$ puan, DA gücünü $+$\%0,7 hatayla yeniden üretmektedir. Verim farkı, modelde \SI{24.0}{\watt}, deneyde \SI{19.7}{\watt} kayba karşılık gelmektedir. Akım THD'si modelde \%39,5, deneyde \%15'tir. Fark dalga biçiminde de görülmektedir. Modelde faz akımı sivri bir tepeden sonra hızlı iner, kısa bir çukurdan sonra az yükselir ve sıfırdan geçer (Şekil~\ref{fig:s1-s2-kalibre-dalga}); osiloskop kaydında tepe yuvarlaktır ve tepeyi bir kambur izler (Şekil~\ref{fig:deney-klasik-faz}). Parametre denemelerinde dalga biçimi IGBT parametrelerinden az, şebeke eşdeğerinin $L$ ve $R$ değerleri ile referans fazından daha çok etkilenmiştir. Aynı değişkenler RMS akımı ve $\cos\varphi$ değerini de değiştirdiğinden kalibrasyon, iki serbest değişkeni güç hedeflerine ayırmıştır.

\textbf{Şebeke gerilimi bozulması.} Benzetim modelinde şebeke ideal sinüs kaynağıdır. Analizör ise varyak ile evirici arasında S1'de \%3,7, S2'de \%5,4, S3'te \%4,8 gerilim THD'si göstermektedir (Tablo~\ref{tab:deney-kosullari}). Ölçüm noktası şebekenin kendisi olmadığından okuma evirici akımının etkisini de içerir. Büyüklük mertebesi için: \SI{120}{\volt}'ta \%3,7 yaklaşık \SI{4.4}{\volt} harmonik gerilimdir ve $L_1+L_2$ toplamının beşinci harmonikteki yaklaşık \SI{10}{\ohm} empedansı üzerinde yaklaşık \SI{0.44}{\ampere} harmonik akım sürebilir; değer \SI{2.3}{\ampere} noktasında temel akımın yaklaşık \%20'sidir. Etki, faz ilişkisine göre evirici harmonikleriyle toplanır ya da onları kısmen söndürür. Etki, düşük akımdaki mutlak THD karşılaştırmasında belirgin bir paya sahiptir.

\umitnot{Jüri sorusuna hazırlık: S1 THD}{%
\textbf{``Model klasik anahtarlamada \%40 veriyor, deney \%15; model yanlış değil mi?''} Model yaklaşık değer verir: verim ve THD sıralaması deneyle aynıdır, S3 THD'si \%14,8'e karşı \%14'tür, S2 dalga biçimi osilogramla aynı karakterdedir. S1 farkı açıkça yazılmıştır; model ideal şebeke kullanır ve $L$--$R$ eşdeğeri ölçülmemiştir.

\textbf{``Aynı ekranda PF 0,921 var, \%15 ile tutmuyor.''} Tezde PF tartışılmıyor ama jüri ekranda görebilir. Ayrıca özet tablodaki $P$, $V$, $I$ aynı ekrandan geldiği için $P/(VI)=0{,}922$ hesabını kendisi yapabilir; \%15 THD ile değer yaklaşık 0,989 olurdu. Cevap: ekran tek anlık okumadır, akım eğrisi kayıt boyunca 2,0--\SI{2.3}{\ampere} arasında değişir, THDi ise süreli kayıttır ve osilogram yaklaşık \%21 verir.}

\textbf{Senaryo sıralaması.} Verim sıralaması S1 $<$ S2 $<$ S3 ve THD sıralaması S3 $<$ S1 $<$ S2, modelde ve deneyde aynıdır. S1$\rightarrow$S3 verim artışı modelde $+3{,}10$, deneyde $+2{,}0$ puan; S1$\rightarrow$S2 artışı modelde $+2{,}78$, deneyde $+1{,}7$ puandır. Fark esas olarak S1 verimindeki $-1{,}35$ puanlık sapmadan gelmektedir.

\textbf{Benzetim tarafı belirsizliği.} Kararlı durum değerleri \SI{0.1}{\second} benzetimin son şebeke periyodundan alınmıştır. S2 ve S3'te son periyot ile 0,08--\SI{0.10}{\second} ortalaması arasındaki fark DA güçte 0,4--\SI{0.5}{\watt}, verimde en çok 0,1 puandır ve verim karşılaştırmasının benzetim tarafı belirsizliğini oluşturur. Kalibre benzetim sonuçları, modelin deneyin büyüklüklerini Tablo~\ref{tab:kalibrasyon-karsilastirma} içinde belirtilen hatayla yeniden ürettiğini göstermektedir.

\section{Bölüm Sonucu}

Deney düzeneği; T-tipi güç katı, faz modülleri, kontrol, tampon ve sürücü kartları, gerilim ve akım ölçüm devreleri, koruma devresi, AA ve DA analizörleri, şönt bağlantısı, PC yazılımı ve üç yönteme ait ekran kayıtlarıyla belgelenmiştir. Deney sonuçları, PC ve osiloskop ekranları ile deney özet tablosu üzerinden raporlanmıştır.

Üç yönteme ait kayıtlar, nötr susturmanın düşük güçte verimi artırırken harmonik bozulmayı büyüttüğü, Sin\textsuperscript{3} teriminin ise susturma kaynaklı bozulmayı azalttığı mekanizmasını desteklemektedir; analizör ve osiloskop kayıtlarının faz akımı THD değerleri birbiriyle uyumludur. Aktif ve reaktif güce kalibre edilen benzetim modeli, S3'te deneyin AA akımını \%1,4, verimini 0,25 puan ve kalibre edilmemiş akım THD'sini 0,8 puan içinde yeniden üretmekte ve deneyin verim ve THD sıralamasını korumaktadır. Özet ve ekran değerleri, her deneyin çalışma koşulları ve ölçüm zinciri tablolar hâlinde verilmiştir. Orta ve anma akımındaki geri beslemeli (S4) çalışma, S4 benzetim taramasıyla incelenmiştir. Faz kolu ölçeğindeki iki katkının kanıt zinciri Bölüm~\ref{chap:deney} ile tamamlanmaktadır; Bölüm~\ref{chap:gttipi}, aynı T-tipi donanımın sistem ölçeğinde şarj ve deşarj işlevlerine yeniden görevlendirilmesini ele almaktadır.
